\documentclass[11pt,a4paper]{article}

\usepackage[margin=2.2cm]{geometry}
\usepackage[T1]{fontenc}
\usepackage{lmodern}
\usepackage{amsmath}
\usepackage{booktabs}
\usepackage{tabularx}
\usepackage{enumitem}
\usepackage{xurl}
\usepackage{hyperref}
\usepackage{xcolor}

\hypersetup{colorlinks=true, linkcolor=blue!50!black, urlcolor=blue!50!black, citecolor=blue!50!black}
\setlist{nosep, leftmargin=1.4em}
\setlength{\parskip}{0.4em}
\setlength{\parindent}{0pt}

\title{Research Gap Discovery: Combining Semantic Embeddings and Citation Structure}
\author{Vedant Acharya (23110010) \and Bhavya Parmar (23110059)}
\date{CS 328 Introduction to Data Science, November 2026}

\begin{document}
\maketitle

\noindent Code: \url{https://github.com/bhavya-parmar/RGD}\\
Data: retrieved from the OpenAlex API (CC0) \cite{priem2022openalex}. The corpus, embeddings and reports are regenerated by the pipeline (Section~\ref{sec:repro}).

\section{Problem Description}

Researchers entering a field must find, among thousands of papers, the places where work is missing. Keyword search returns papers that share vocabulary, and citation counts favour older, well-known work. Neither shows \emph{research silos}: groups of papers that work on closely related problems but rarely cite each other. Such silos are candidate research gaps, because results from one group may transfer to the other.

We define the task as follows. Given a topic query, build a corpus of related papers and find pairs of paper communities with
\begin{itemize}
  \item high semantic similarity (they study related problems), and
  \item low citation connectivity (they rarely cite each other),
\end{itemize}
and test whether each such pattern is statistically supported rather than an artefact of chance. We also report individual paper pairs that are highly similar but have no citation link in either direction.

The case study in this report uses the query \emph{anisotropy in large language models}: the observation that contextual embeddings occupy a narrow cone of the vector space.

\section{Solution Approach}

\subsection{Survey of existing methods}

Citation-based literature mapping tools.
\begin{itemize}
  \item Connected Papers, Litmaps, ResearchRabbit and Inciteful build graphs around one or more seed papers using co-citation, bibliographic coupling or direct citation \cite{tay2022tools}. Inciteful ranks papers with PageRank and Adamic-Adar similarity and can find paths between two papers. These tools help explore the neighbourhood of known papers but do not compare citation structure with content, and do not test whether a pattern is significant.
  \item VOSviewer \cite{vaneck2010vosviewer} and CiteSpace \cite{chen2006citespace} are bibliometric science-mapping tools that cluster co-citation and co-occurrence networks. They describe the structure of a field but do not flag pairs of clusters that should be connected and are not.
\end{itemize}

Comparing citation and text views of science.
\begin{itemize}
  \item Boyack and Klavans compared co-citation, bibliographic coupling, direct citation and text-based relatedness on millions of biomedical papers \cite{boyack2010citation}, and found that citation-text hybrids give the most accurate clusters \cite{boyack2020hybrid}. Their aim is better clustering, not detecting disagreement between the two views.
  \item Recent work uses embeddings to add missing edges to fragmented citation graphs \cite{semaug2026}, which treats a semantic link without a citation as missing data. We treat the same disagreement as the signal of interest.
\end{itemize}

Literature-based discovery (LBD).
\begin{itemize}
  \item Swanson's ABC model links two disjoint literatures A and C through shared intermediate concepts B, as in the fish oil and Raynaud's disease discovery \cite{swanson1986fish}. Modern LBD uses knowledge graphs and embeddings to predict such links \cite{lbdsurvey2025}. LBD works at the level of concepts, usually in biomedicine, whereas we work at the level of papers and citation communities in any field.
\end{itemize}

Building blocks.
\begin{itemize}
  \item SPECTER \cite{cohan2020specter} and SPECTER2 \cite{singh2023scirepeval} embed a paper's title and abstract with a transformer trained on citation triplets, so papers that cite each other are close in embedding space.
  \item MinHash \cite{broder1997resemblance} with locality-sensitive hashing (LSH) finds near-duplicate documents without comparing all pairs; FAISS \cite{johnson2019faiss} does the same for dense vectors.
  \item Louvain \cite{blondel2008louvain} detects communities by maximising modularity; k-means++ \cite{arthur2007kmeanspp} clusters vectors.
  \item The Benjamini-Hochberg procedure \cite{benjamini1995fdr} controls the false discovery rate when many hypotheses are tested.
\end{itemize}

\subsection{Our approach}

The pipeline has four phases. Configuration is in \texttt{config.yaml}; default values are given below.

Phase 1: retrieval and deduplication.
\begin{itemize}
  \item Query OpenAlex for seed papers. An optional \texttt{--must-include} flag adds a boolean filter on title and abstract (for example \texttt{anisotropy AND ("language model" OR transformer)}), because plain OpenAlex search matches full text and ranks loosely.
  \item Expand the corpus with each seed's references (up to 20) and OpenAlex related works.
  \item Remove duplicates: exact match on ID and normalised title, then MinHash LSH on word 3-shingles of title and abstract (128 hash functions). LSH candidates are verified with exact Jaccard similarity $J \geq 0.8$, and the LSH threshold is set 0.15 lower so that true duplicates are rarely missed.
\end{itemize}

Phase 2: representation.
\begin{itemize}
  \item Embed \texttt{title [SEP] abstract} with SPECTER2 (\texttt{[CLS]} token, 768 dimensions). The proximity adapter is used when the installed library versions allow it; otherwise the base model.
  \item Reduce with PCA to $\min(64, n-1)$ dimensions and L2-normalise. PCA centres the data, which removes the direction shared by all embeddings (Section~\ref{sec:results}).
  \item Index the vectors in FAISS (exact inner product, equal to cosine similarity).
  \item Build a directed citation graph with an edge $u \to v$ when corpus paper $u$ cites corpus paper $v$.
\end{itemize}

Phase 3: communities and frontier scoring.
\begin{itemize}
  \item Detect citation communities (silos) with Louvain on the undirected graph. Communities with fewer than 3 papers are ignored.
  \item For each pair of communities $(A, B)$ compute
  \begin{align*}
    \mathrm{sim}(A,B) &= \frac{1}{|A||B|}\sum_{a \in A}\sum_{b \in B} a^\top b = \bar{a}^\top \bar{b}, \\
    \mathrm{density}(A,B) &= \frac{\#\{\text{citations } A \to B \text{ and } B \to A\}}{2|A||B|}, \qquad
    \mathrm{ratio}(A,B) = \frac{\mathrm{density}(A,B)}{\mathrm{density}(G)}, \\
    \mathrm{score}(A,B) &= \mathrm{sim}(A,B)\,\bigl(1 - \mathrm{density}(A,B)\bigr).
  \end{align*}
  The identity $\mathrm{sim} = \bar{a}^\top\bar{b}$ replaces an $|A| \times |B|$ matrix with one dot product. A pair is low-citation if $\mathrm{ratio} \leq 1$, that is, its members cite each other less than an average pair of papers in the corpus.
  \item Isolated paper pairs: FAISS range search returns all pairs above the 95th percentile of pairwise similarity; pairs with a citation in either direction are removed. The percentile adapts the cutoff to the embedding space.
\end{itemize}

Phase 4: validation and report.
\begin{itemize}
  \item Similarity confidence interval: resample 80\% of papers with replacement 500 times, recompute $\mathrm{sim}(A,B)$ for the fixed communities, and report the 2.5th and 97.5th percentiles.
  \item Permutation test. $H_0$: $A$ and $B$ are no more similar in topic than two random groups of the same sizes. Shuffle which embedding belongs to which paper (communities and citations fixed), recompute the score, and set $p = (\#\{\text{null} \geq \text{observed}\} + 1)/(N + 1)$. Groups come from citations and the null from embeddings, so the test is not circular.
  \item All community pairs are tested; p-values are adjusted with Benjamini-Hochberg and only then truncated to the top 10. A pair is reported as a frontier if it is significant at $\alpha = 0.05$ after adjustment and is low-citation.
  \item The HTML report labels each community by its most typical paper: the member closest to the community's mean embedding.
\end{itemize}

\section{Results}
\label{sec:results}

Run settings: query \emph{anisotropy in large language models}, default configuration, 100 permutations.

\subsection{Corpus and representation}

\begin{table}[h]
\centering
\caption{Corpus, embedding and graph statistics.}
\label{tab:stats}
\begin{tabular}{lr}
\toprule
Records retrieved (seeds and citation expansion) & 196 \\
Papers after deduplication & 195 \\
Papers mentioning anisotropy in title or abstract & 95 \\
Publication years (median) & 1999 to 2026 (2024) \\
\midrule
Raw SPECTER2 cosine: mean (5th, 95th percentile) & 0.879 (0.802, 0.935) \\
After PCA (64 comp., 90.9\% var.): mean (95th pct.) & $-0.002$ (0.350) \\
k-means ($k=8$) mean similarity within / between clusters & 0.237 / $-0.037$ \\
\midrule
Citation edges between corpus papers & 429 \\
Graph density & 0.0113 \\
Papers with no citation link to the corpus & 60 \\
Louvain communities (modularity) & 68 (0.390) \\
Communities with at least 3 papers (sizes) & 7 (25, 25, 22, 20, 15, 14, 12) \\
Community pairs tested & 21 \\
\bottomrule
\end{tabular}
\end{table}

The raw SPECTER2 embeddings are themselves strongly anisotropic: two unrelated papers have cosine similarity around 0.8. Fixed similarity thresholds are therefore meaningless on raw vectors. PCA centring removes the common direction and spreads similarities around 0, which is why all cutoffs are percentiles.

\subsection{Community frontiers}

Of 21 community pairs, 14 are low-citation and 5 are significant after correction. No pair is both, so the run reports no frontier. Table~\ref{tab:frontiers} lists the significant pairs and the closest low-citation candidates. Communities are labelled by their most typical paper.

\begin{table}[h]
\centering
\caption{Selected community pairs. Ratio: citation density relative to the whole graph. $p_{\mathrm{adj}}$: Benjamini-Hochberg adjusted p-value (100 permutations; for the last two rows, 5000 permutations).}
\label{tab:frontiers}
\small
\begin{tabularx}{\textwidth}{Xrrrrr}
\toprule
Community pair & sim [95\% CI] & cross-cit. & ratio & $p_{\mathrm{adj}}$ & frontier \\
\midrule
Sentence embeddings (Ditto) / Representation degeneration & 0.085 [0.035, 0.149] & 35 & 4.12 & 0.042 & no \\
LERT (pre-trained LMs) / Sentence embeddings (Ditto) & 0.058 [0.008, 0.115] & 22 & 1.55 & 0.042 & no \\
Sentence embeddings (Ditto) / BERTology & 0.050 [$-0.010$, 0.117] & 36 & 3.18 & 0.042 & no \\
LERT / Representation degeneration & 0.039 [$-0.013$, 0.091] & 12 & 1.41 & 0.042 & no \\
LERT / Pythia (training dynamics) & 0.024 [$-0.038$, 0.085] & 8 & 1.18 & 0.042 & no \\
\midrule
SAE steering of LLMs / Pythia & 0.028 [$-0.056$, 0.110] & 3 & 0.50 & 0.057 & no \\
Pythia / BERTology & 0.026 [$-0.030$, 0.077] & 5 & 0.92 & 0.069 & no \\
\bottomrule
\end{tabularx}
\end{table}

Observations.
\begin{itemize}
  \item Every pair with significant topic overlap cites more than average (ratio 1.18 to 4.12). In this field, citation silos that share a topic already cite each other.
  \item The strongest candidate, sparse-autoencoder steering of LLMs versus Pythia-style training-dynamics analysis, cites at half the average rate but misses significance ($p_{\mathrm{adj}} = 0.057$). The result is unchanged with 1000 and 5000 permutations, so it is not limited by permutation resolution.
  \item Several confidence intervals include 0, so mean similarities between communities are small and uncertain at this corpus size.
  \item On an earlier 145-paper corpus for the same query, one pair passed ($p_{\mathrm{adj}} = 0.049$). It did not reappear when the corpus grew to 195 papers, so borderline frontiers should be treated as leads, not findings.
\end{itemize}

\subsection{Isolated paper pairs}

The paper-level output is the most directly useful. The 95th percentile cutoff was 0.350, and 15 pairs were reported. Relevant examples with no citation in either direction:
\begin{itemize}
  \item IsoScore: Measuring the Uniformity of Embedding Space Utilization / Frequency-based Distortions in Contextualized Word Embeddings (0.710).
  \item Is anisotropy really the cause of BERT embeddings not being semantic? / Frequency-based Distortions in Contextualized Word Embeddings (0.703).
  \item Activation Scaling for Steering and Interpreting Language Models / Steering Large Language Model Activations in Sparse Spaces (0.699).
  \item Ditto: A Simple and Efficient Approach to Improve Sentence Embeddings / WhiteningBERT: An Easy Unsupervised Sentence Embedding Approach (0.778).
\end{itemize}
The top-ranked pair (0.914) and three others are cosmology papers on \emph{Lorentz anisotropy}, retrieved because the query term is ambiguous. They are similar to each other and absent from the citation graph, so they pass both criteria. This shows that corpus filtering directly affects the output.

\section{Conclusions}

\begin{itemize}
  \item We built an end-to-end pipeline that retrieves a corpus from OpenAlex, deduplicates it with verified MinHash LSH, represents it with SPECTER2 embeddings and a citation graph, scores citation-community pairs for high similarity and low citation, and validates them with bootstrap confidence intervals and a Benjamini-Hochberg corrected permutation test.
  \item For the anisotropy corpus, no community pair is both significantly similar and under-cited. Topic overlap and citation are positively related in this field. The method reports this as a negative result instead of ranking noise as gaps.
  \item Paper-level isolated pairs give concrete reading suggestions, such as connecting work on uniformity metrics (IsoScore) with work on frequency-based distortions.
  \item Design lessons: raw scientific embeddings are anisotropic and need centring before any similarity threshold; comparing different k-means clusters cannot find "similar but separate" groups; a hypothesis test must draw its null from a different view than the one that defines the groups; and when many pairs are tested, multiple-testing correction is necessary.
  \item Limitations: the corpus is small (195 papers, 429 internal citations), and 60 papers have no citation link to the rest; OpenAlex citation coverage and abstract availability are incomplete; ambiguous query terms pull in off-topic papers; Louvain yields many singleton communities on sparse graphs; and the bootstrap measures the stability of similarity for fixed communities, not whether the same frontier would be detected again.
  \item Future work: larger corpora with deeper citation expansion, bibliographic coupling or co-citation as additional citation signals, domain filtering of the corpus, and a bootstrap that re-runs community detection on each sample.
\end{itemize}

\section{Code, data and reproduction}
\label{sec:repro}

\begin{itemize}
  \item Code: \url{https://github.com/bhavya-parmar/RGD} (pipeline in \texttt{src/}, settings in \texttt{config.yaml}, method notes in \texttt{docs/}).
  \item Data: OpenAlex works API (\url{https://openalex.org}), CC0 licensed. A free API key in \texttt{.env} is required.
  \item Reproduction:
\end{itemize}
\begin{verbatim}
pip install -r requirements.txt
python src/pipeline.py --query "anisotropy in large language models" \
    --must-include anisotropy "language model|transformer" --max-seeds 50
\end{verbatim}
Outputs are written to \texttt{reports/} (HTML report, validated frontiers and isolated paper pairs as CSV), with the intermediate corpus, embeddings and FAISS index under \texttt{data/} and \texttt{models/}.

\begin{thebibliography}{99}
\bibitem{priem2022openalex} J. Priem, H. Piwowar, R. Orr. OpenAlex: A fully-open index of scholarly works, authors, venues, institutions, and concepts. arXiv:2205.01833, 2022.
\bibitem{tay2022tools} A. Tay. Citation-based literature mapping tools: a comparison. 2022. \url{https://aarontay.substack.com/p/citation-based-literature-mapping-tools}
\bibitem{vaneck2010vosviewer} N. J. van Eck, L. Waltman. Software survey: VOSviewer, a computer program for bibliometric mapping. Scientometrics 84(2), 2010.
\bibitem{chen2006citespace} C. Chen. CiteSpace II: Detecting and visualizing emerging trends and transient patterns in scientific literature. JASIST 57(3), 2006.
\bibitem{boyack2010citation} K. W. Boyack, R. Klavans. Co-citation analysis, bibliographic coupling, and direct citation: Which citation approach represents the research front most accurately? JASIST 61(12), 2010.
\bibitem{boyack2020hybrid} K. W. Boyack, R. Klavans. A comparison of large-scale science models based on textual, direct citation and hybrid relatedness. Quantitative Science Studies 1(4), 2020.
\bibitem{semaug2026} Reconnecting fragmented citation networks with semantic augmentation. arXiv:2605.12263, 2026.
\bibitem{swanson1986fish} D. R. Swanson. Fish oil, Raynaud's syndrome, and undiscovered public knowledge. Perspectives in Biology and Medicine 30(1), 1986.
\bibitem{lbdsurvey2025} Recent advances and future directions in literature-based discovery. arXiv:2506.12385, 2025.
\bibitem{cohan2020specter} A. Cohan, S. Feldman, I. Beltagy, D. Downey, D. S. Weld. SPECTER: Document-level representation learning using citation-informed transformers. ACL, 2020.
\bibitem{singh2023scirepeval} A. Singh, M. D'Arcy, A. Cohan, D. Downey, S. Feldman. SciRepEval: A multi-format benchmark for scientific document representations. EMNLP, 2023 (arXiv:2211.13308).
\bibitem{broder1997resemblance} A. Z. Broder. On the resemblance and containment of documents. Compression and Complexity of Sequences, 1997.
\bibitem{johnson2019faiss} J. Johnson, M. Douze, H. J\'egou. Billion-scale similarity search with GPUs. IEEE Transactions on Big Data, 2019.
\bibitem{blondel2008louvain} V. D. Blondel, J.-L. Guillaume, R. Lambiotte, E. Lefebvre. Fast unfolding of communities in large networks. J. Stat. Mech., 2008.
\bibitem{arthur2007kmeanspp} D. Arthur, S. Vassilvitskii. k-means++: The advantages of careful seeding. SODA, 2007.
\bibitem{benjamini1995fdr} Y. Benjamini, Y. Hochberg. Controlling the false discovery rate: a practical and powerful approach to multiple testing. J. R. Stat. Soc. B 57(1), 1995.
\end{thebibliography}

\end{document}
